%% ==========================================================================
%% SECTION 1: INTRODUCTION
%% ==========================================================================
\balance
\section{Introduction}\label{sec:intro}

Systems incorporating language model components face a verification problem.
The components generate text by sampling from learned probability distributions,
and some fraction of that text is fabricated: fluent, confident, and wrong.
LLM-based systems today are rife with ``hallucinations''.

% This paper provides the formal framework and empirical cost surface for making that decision.

We use the term \emph{fabrication} throughout this paper to emphasize that
the system's default behavior under coherence optimization is reasonable and
not a malfunction to be patched. Language models optimized for
coherence produce coherent output. The output reflects the training data: grounded
training data tends to track truth, while ungrounded training data tends to produce untruth.
\citeauthor{xu2024hallucination}~\cite{xu2024hallucination} proves that hallucination is
inevitable for LLMs over the class of computable functions.
% Our result is orthogonal. They prove a computational limit on what the model can produce;
% we prove an observational limit on what a supervisor can verify.
% [moved to telemetry paragraph below]

%Verification has a cost. 
Consequently, every system that
deploys a language model component must allocate a budget: how much of the output to verify,
which outputs to prioritize, and what signals to use for triage.
For example, checking a citation against a database costs API
calls and latency; having a domain expert review a summary costs time
and salary; running a second model as a judge costs compute.
Verification cost varies by content type: checking a citation is cheap;
assessing whether a summary faithfully represents its source is expensive.

The key issue is that the current text-only interface to language models does not expose the model's
internal computation, which produces signals that distinguish grounded generation from fabrication.
The model may know it is fabricating but currently the supervisor cannot tell if the model
is fabricating just from the text output. 

This observability gap has direct consequences for verification cost. 
A supervisor receiving only text must spend its
verification budget reconstructing, from the output alone,
what the model's own computation already knew. The budget is
consumed by reconstruction rather than concentrated on the cases
where verification is most needed.

Our key insight is that language models carry internal
signals, which are the entropy traces, attention geometry, token-level
log-probabilities, that distinguish grounded generation from
fabrication. These signals are byproducts of the computation itself:
\emph{telemetry} that the model cannot separately control under standard training.
Unlike the model's \emph{testimony} (the text it chooses to output),
telemetry is harder to fake.

We introduce \emph{epistemic observability}: an interface that exposes this
telemetry alongside the standard text output.
Unlike interpretability, which seeks to explain why a model activates particular neurons or circuits, 
or explainability, which generates post-hoc rationales for its outputs, 
epistemic observability addresses a different question: how trustworthy is this generation? Rather than explaining the model's internal reasoning, it surfaces signals produced as a byproduct of inference that help estimate whether the generated content is grounded in the model's knowledge or is likely to be fabricated.

To support a range of application requirements, we propose a multi-tier epistemic observability interface, 
where each successive tier exposes richer and more informative telemetry about the generation process. 
These tiers provide a principled trade-off between observability and computational cost, 
enabling developers to select the level of epistemic insight that best matches their system's reliability requirements and resource budget. 
For instance, a creative brainstorming assistant may require only Tier~0 observability, 
whereas a medical decision-support system may justify the additional cost of higher observability tiers 
to better identify and validate potentially hallucinated content. In this way, the observability hierarchy 
makes explicit what additional epistemic information each level of investment provides.

As a proof of concept, we build a new tensor interface at Tier 1 epistemic observability,
which includes the entropy at the token granularity. \todo{1 liner about the eval}

% \fakepara{Empirical cost surface of verification.} We map the cost
% surface across text-only and tensor-augmented verification strategies.
% Under length-controlled evaluation, no text-channel feature exceeds pooled
% AUC 0.70; per-token entropy is the only signal that generalizes across
% architectures. Tensor-guided triage outperforms text baselines at every
% budget level (\autoref{sec:eval}).
